% !TeX root = ../main.tex

\section{Method}\label{section:design}

In this section, we present our three-stage personalized adaptation framework, illustrated in Figure~\ref{fig:method-overview}. We begin with a formal problem formulation, then describe each component in detail: doctor-specific LoRA training via LLaMA-Factory~\cite{zheng2024llamafactory}, Doctor-Adaptive Routing (DAR) for dynamic adapter routing, and SCF for parameter fusion.

\subsection{Problem Formulation}

Let $\mathcal{D} = \{d_1, d_2, \ldots, d_K\}$ denote the set of $K$ physicians in an internet hospital, and let $\mathcal{M}_0$ denote a pre-trained base language model. Each physician $d_k$ is associated with a historical dialogue corpus $\mathcal{C}_k = \{(q_i^{(k)}, r_i^{(k)})\}_{i=1}^{N_k}$, where $q_i^{(k)}$ is a patient query and $r_i^{(k)}$ is the physician's corresponding response. Our objective is to construct a system that, given a new patient query $q$ directed to physician $d_k$, generates a candidate response $\hat{r}$ that (1) is clinically accurate, (2) reflects physician $d_k$'s personal linguistic style and departmental norms, and (3) complies with institutional regulations.

Formally, we seek to learn physician-specific adaptation parameters $\Delta \Theta_k$ such that:
\begin{equation}
\hat{r} = \mathcal{M}_{\Theta_0 + \Delta \Theta_k}(q), \quad \text{where } \Delta \Theta_k \text{ captures the personalization for doctor } d_k.
\end{equation}

The key challenges are: (a) $N_k$ is typically small (tens to hundreds of examples), making full fine-tuning infeasible; (b) $K$ may be large, requiring efficient multi-doctor serving; and (c) the personalization signal must be disentangled from general clinical competence encoded in $\mathcal{M}_0$.

% % Method overview flowchart
\begin{CJK}{UTF8}{gbsn}
\begin{figure}[H]
\centering
\begin{tikzpicture}[
    node distance=0.65cm,
    box/.style={draw, rounded corners=3pt, minimum width=3.2cm, minimum height=0.7cm, align=center, font=\footnotesize},
    stagebox/.style={draw, rounded corners=3pt, fill=blue!10, draw=blue!50, minimum width=3.5cm, minimum height=0.75cm, align=center, font=\footnotesize\bfseries},
    inputbox/.style={draw, rounded corners=3pt, fill=green!8, draw=green!50!black, minimum width=3.5cm, minimum height=0.6cm, align=center, font=\footnotesize},
    outputbox/.style={draw, rounded corners=3pt, fill=orange!10, draw=orange!60!black, minimum width=5cm, minimum height=0.65cm, align=center, font=\footnotesize},
    arrow/.style={thick, ->, >=stealth},
    dasharrow/.style={thick, ->, >=stealth, dashed, blue!60},
]

% ===== STAGE 1: LoRA Training (x=0) =====
\node[stagebox] (s1) at (0,0) {Stage 1: 医生专属LoRA训练};
\node[inputbox] (s1-input) at (0,1.8) {医生历史问诊数据\\(IHDPD 数据集)};
\node[box] (s1-weclone) at (0,-1.3) {weclone 训练流水线};
\node[box] (s1-lora) at (0,-2.4) {LoRA 适配器 $\Delta\Theta_1,\ldots,\Delta\Theta_K$};
\node[font=\footnotesize\itshape, text=black!50] at (0,-3.0) {每个适配器 $<$10 MB};

\draw[arrow] (s1-input) -- (s1);
\draw[arrow] (s1) -- (s1-weclone);
\draw[arrow] (s1-weclone) -- (s1-lora);

% ===== STAGE 2: LAG Routing (x=4.8) =====
\node[stagebox] (s2) at (4.8,0) {Stage 2: LAG 智能路由};
\node[inputbox] (s2-input) at (4.8,1.8) {患者问询 $q$};
\node[box] (s2-id) at (4.8,-1.3) {身份路由~(目标医生已知)};
\node[box] (s2-qr) at (4.8,-2.15) {查询路由~(自动匹配风格)};
\node[font=\footnotesize\itshape, text=black!50] at (7.5,-2.15) {Sentence-BERT 编码};

\draw[arrow] (s2-input) -- (s2);
\draw[arrow] (s2.south) -- (s2-id.north);
\draw[arrow] (s2-id.south) -- (s2-qr.north);

% Arrow: LoRAs -> LAG with label node
\node[font=\footnotesize, align=center] (lora2lag) at (2.4,-1.8) {按需加载\\对应适配器};
\draw[arrow] (1.6,-2.4) to[bend left=10] (3.2,-2.4);

% ===== STAGE 3: TIES-MERGING (x=9.6) =====
\node[stagebox] (s3) at (9.6,0) {Stage 3: TIES-MERGING};
\node[box] (s3-trim) at (9.6,-1.1) {TRIM: 裁剪低幅值参数};
\node[box] (s3-sign) at (9.6,-1.95) {ELECT SIGN: 符号冲突消解};
\node[box] (s3-merge) at (9.6,-2.8) {MERGE: 加权平均融合};
\node[box, fill=blue!8] (s3-result) at (9.6,-3.65) {统一融合适配器 $\Delta\Theta_{\mathrm{merged}}$};

\draw[arrow] (s3) -- (s3-trim);
\draw[arrow] (s3-trim) -- (s3-sign);
\draw[arrow] (s3-sign) -- (s3-merge);
\draw[arrow] (s3-merge) -- (s3-result);

% Dashed arrow: all LoRAs -> TIES
\node[font=\footnotesize, align=center] (lora2ties) at (5.5,-3.8) {多适配器\\离线融合};
\draw[dasharrow] (1.6,-3.3) to[bend right=20] (8.0,-3.3);

% ===== Output =====
\node[outputbox] (output) at (4.8,-5.5) {\textbf{个性化医生响应} $\hat{r}$\\\footnotesize 贴合目标医生的语言风格、科室规范和临床习惯};

% Arrows to output
\draw[arrow] (s2-qr.south) -- ++(0,-0.7) -| (output.north);
\draw[arrow] (s3-result.south) -- ++(0,-0.7) -| (output.north);

% ===== Legend =====
\node[font=\footnotesize] at (4.8,-6.3) {Stage 1 (离线训练)~~$\longrightarrow$~~Stage 2 (在线推理)~~$\longrightarrow$~~Stage 3 (离线融合)~~$\longrightarrow$~~多医生单端点部署};

\end{tikzpicture}
\caption{三阶段个性化适配框架整体流程。Stage~1：利用 weclone 流水线为每位医生训练独立的 LoRA 适配器；Stage~2：通过 LAG 路由根据患者问询和目标医生身份动态选择适配器；Stage~3：使用 TIES-MERGING (TRIM $\rightarrow$ ELECT SIGN $\rightarrow$ MERGE) 将多个医生适配器融合为统一部署工件，支持单端点服务所有医生。}
\label{fig:method-overview}
\end{figure}
\end{CJK}
\begin{figure}[!t]
\includegraphics[width=\textwidth]{figures/design.png}
\caption{Overview of the three-stage personalized adaptation framework. Offline: Doctor-specific LoRA adapters are trained offline via LoRA fine-tuning using each physician's historical consultation records. Online: At inference time, the DAR router dynamically selects the appropriate adapter based on the patient query and target physician identity. SCF fuses multiple doctor-specific adapters into a single unified deployment artifact with LoRAX, enabling efficient multi-physician serving through a single inference endpoint.}
\label{fig:method-overview}
\end{figure}   
\unskip


\subsection{Doctor-Specific LoRA Training}

We adopt Low-Rank Adaptation (LoRA)~\cite{hu2021lora} as our parameter-efficient fine-tuning strategy. For each physician $d_k$, we train a LoRA adapter that modifies the attention projection matrices of $\mathcal{M}_0$. Specifically, for a pre-trained weight matrix $W \in \mathbb{R}^{d \times d}$, LoRA constrains its update to a low-rank decomposition:
\begin{equation}
W + \Delta W = W + B A,
\end{equation}
where $B \in \mathbb{R}^{d \times r}$ and $A \in \mathbb{R}^{r \times d}$ are trainable low-rank matrices with rank $r \ll d$. During training, only $A$ and $B$ are optimized, while $W$ remains frozen.

We adopt LoRA fine-tuning based on LLaMA-Factory~\cite{zheng2024llamafactory} (via WeClone~\cite{xmingXming521WeClone2026}), an efficient framework for parameter-efficient fine-tuning that provides optimized kernels for training and inference. For each physician $d_k$, the training objective is the standard causal language modeling loss:
\begin{equation}
\mathcal{L}_k = -\sum_{t=1}^{T} \log P(r_t^{(k)} \mid q^{(k)}, r_{<t}^{(k)}; \Theta_0 + \Delta \Theta_k).
\end{equation}

To prevent overfitting given the limited per-physician data, we apply the following regularization strategies: (1) a small LoRA rank $r \in \{4, 8, 16\}$; (2) dropout on the LoRA matrices with rate $\alpha = 0.1$; and (3) early stopping based on validation perplexity on a held-out subset of each physician's dialogues. Each trained adapter $\Delta \Theta_k$ is serialized to disk as a compact checkpoint (typically $<10$ MB), enabling efficient storage and loading.

\subsection{Doctor-Adaptive Routing (DAR)}

At inference time, the system must select the appropriate LoRA adapter for the target physician. Let $e_k \in \mathbb{R}^{384}$ denote a pre-computed embedding vector encoding physician $d_k$'s communication style profile (derived from their historical responses via Sentence-BERT~\cite{reimers2019sentence}). For a new patient query $q$, we compute its embedding $e_q$ using the same encoder. The DAR router defines a scoring function:

\begin{equation}
\text{score}(q, d_k) = \cos(e_q, e_k) = \frac{e_q^\top e_k}{\|e_q\| \cdot \|e_k\|}
\end{equation}

We support two routing modes:

\textbf{Identity-Based Routing.} When the target physician $d_{k^*}$ is known at query time (the standard internet hospital scenario), the router directly selects adapter $k^*$ with zero computational overhead:
\begin{equation}
k_{\text{selected}} = k^*
\end{equation}

\textbf{Query-Based Routing.} When the optimal physician is not predetermined (e.g., triage or automated preliminary response generation), the router retrieves the top-$m$ most similar physicians by cosine similarity to the query embedding:
\begin{equation}
\mathcal{K}_m(q) = \underset{k \in \{1,\ldots,K\}}{\text{top-}m}\; \cos(e_q, e_k)
\end{equation}
The selected adapter(s) are then applied to the base model for generation. In practice, the FAISS index~\cite{johnson2019faiss} enables sub-millisecond nearest-neighbor retrieval over the full set of $K$ physician embeddings.

\textbf{Routing Context Augmentation.} Beyond the physician identity and query-text signals described above, the DAR router integrates four additional knowledge sources to enrich the routing context:
\begin{enumerate}
\item \textbf{Department-level regulations.} Each clinical department maintains codified rules governing permissible response content, specialty-specific disclaimers, and mandatory referral triggers. These regulations are stored as structured text snippets and retrieved by department ID at routing time.
\item \textbf{Hospital-wide institutional policies.} Global policies---including patient privacy disclaimers, escalation protocols, and medication prescription guidelines---are maintained as a centralized document set and prepended to the routing context regardless of department.
\item \textbf{Physician--patient dialogue history.} For returning patients, the most recent $N$ turns of the doctor--patient conversation are retrieved to provide continuity of care context, ensuring the generated response is clinically consistent with prior interactions.
\item \textbf{Physician-curated knowledge base.} Each physician maintains a personal editable archive of clinical notes, treatment templates, frequently-used drug combinations, and specialty references. These documents are indexed via a lightweight retrieval-augmented generation (RAG) pipeline built on FAISS~\cite{johnson2019faiss} with Sentence-BERT embeddings. At inference time, the top-$t$ most relevant knowledge snippets are retrieved and injected into the routing prompt, allowing physicians to dynamically refine their knowledge base without retraining the underlying LoRA adapter. Updates to the knowledge base propagate to the RAG index in real time.
\end{enumerate}
The augmented routing context $C_k$ for physician $d_k$ and query $q$ is assembled as: $C_k = \{ \text{dept\_rules}, \text{hospital\_policy}, \text{history}_{N}, \text{RAG\_top}_{t} \}$. This context is prepended to the prompt template before generation, ensuring all regulatory, institutional, and personal knowledge constraints are satisfied simultaneously.

\subsection{SCF for Parameter Fusion}

While per-physician LoRA adapters provide maximal personalization, serving $K$ separate adapters incurs linear storage and loading costs. SCF fuses multiple adapters into a single merged artifact while minimizing interference~\cite{yadavTIESMergingResolvingInterference2023a}. Formally, let $\Delta_k \in \mathbb{R}^{d}$ denote the flattened LoRA parameter delta for physician $d_k$. SCF proceeds in three stages:

\textbf{Stage 1: Trim.} Low-magnitude parameter changes likely representing noise are reset to zero. For each parameter dimension $j$, we retain only the top $(1-\tau)$ fraction of parameters by magnitude:
\begin{equation}
\Delta_{k,j}^{(1)} = \begin{cases}
\Delta_{k,j}, & \text{if } |\Delta_{k,j}| \ge \text{Q}_{\tau}(\{|\Delta_{1,j}|,\ldots,|\Delta_{K,j}|\}) \\
0, & \text{otherwise}
\end{cases}
\end{equation}

\textbf{Stage 2: Elect Sign.} For each dimension, we resolve sign conflicts via a weighted majority vote with soft sign parameter $\eta \ge 0$:
\begin{equation}
s_j^* = \operatorname{sign}\!\left(\sum_{k=1}^{K} |\Delta_{k,j}^{(1)}|^{\eta} \cdot \operatorname{sign}(\Delta_{k,j}^{(1)})\right)
\end{equation}
When $\eta=1$, all deltas contribute equally; as $\eta \to \infty$, only the largest-magnitude delta dominates.

\textbf{Stage 3: Merge.} Only parameters whose sign agrees with $s_j^*$ are merged. An adaptive trim coefficient $\lambda \in [0,1]$ controls merging aggressiveness:
\begin{equation}
\Delta_j^{\text{(merged)}} = \frac{1}{|\mathcal{A}_j|} \sum_{k \in \mathcal{A}_j} \Delta_{k,j}^{(1)},\quad
\mathcal{A}_j = \{k : \operatorname{sign}(\Delta_{k,j}^{(1)}) = s_j^*\}
\end{equation}

The merged adapter is applied as $\Theta^{\text{(merged)}} = \Theta_0 + \Delta^{\text{(merged)}}$. Default hyperparameters are $\eta=1.0$, $\lambda=0.5$, and $\tau=0.20$, with sensitivity analyzed in Section~\ref{section:evaluation}.

